\scanpage{37}
% L1–6, continuation
d'où \inserted{(remplaçant $A$ par $A/J$)} \unsure{(i)|(ii)} $+$ le fait que $M$ est Coh. Mac. Or de ceci on déduit facilement que $A$ est Coh. Mac. et $M\simeq\Omega_A$. On peut essayer une voie analogue pour (iv), en y déduisant de la \uline{première} hyp. \inserted{$+$ la condition du support} que l'annulateur de $\mathfrak m$ dans $M$ \uline{et} $A$ \deleted{sont} nuls, d'où par récurrence que $M$ \uline{et} $A$ sont Coh. Mac, \deleted{\unsure{ce}} qui ramène comme précédemment à un cas du (i)

% L7–13
Démonstration des équivalences dans (i)

$b\Leftrightarrow b'\qquad c\unsure{\Rightarrow|\Leftrightarrow}c'$, \qquad $a\Rightarrow b$ \deleted{\unsure{$\Rightarrow$}} triviaux

$b\Rightarrow c$ car $b\Leftrightarrow b'$, donc $D(H^n(M))\simeq\widehat A$ et $\lg_{\mathfrak p}(A)=\lg_{\mathfrak p}(M)$

$c'\Rightarrow b'$ car $c'\simeq\quad D(H^n(M))\simeq\widehat A^{m}$ et $m\lg_{\mathfrak p}(A)=\lg_{\mathfrak p}(M)$ d'où $m=1$

\margin{\deleted{\unsure{$b\Leftrightarrow b'$|$b'$}} \unsure{$c\Rightarrow c'$} \unsure{équiv.}}

$c\Rightarrow d$, car $N\to\operatorname{Ext}^{n+1}(N,M)$ est exact à gauche, d'où $\operatorname{Ext}^{n+1}(k,M)\to H^{n+1}(M)$ \uline{inj}.

$d\Rightarrow c'$ trivial \deleted{\unsure{donc}} // d'où $a\Rightarrow b\Leftrightarrow b'\Leftrightarrow c\Leftrightarrow c'\Leftrightarrow d$

% L14–16
\[
\begin{aligned}
b&\Rightarrow\operatorname{Ext}^n(k,M)\simeq k,\quad
  \text{et b')}\Rightarrow H^n(M)\text{ est fidèle donc }M\text{ est fidèle, d'où}\\
 &\phantom{\Rightarrow}\operatorname{Ext}^n(A/\unsure{\mathfrak n|\mathfrak N},M)\simeq D(A/\mathfrak n)
\end{aligned}
\]
$(b,b')\Rightarrow e$ et $f$, donc $(b,b',d)\Rightarrow a$

% L17–21
Reste à prouver que f) $\Rightarrow$ a). \inserted{Si $M$ est Coh. Mac.} Or
\[
\begin{aligned}
\operatorname{Ext}^n(A/\mathfrak n,M)&=\operatorname{Hom}(A/\mathfrak n,M)\\
\operatorname{Ext}^n(k,M)&=\operatorname{Hom}(k,\unsure{M|\mu})\\
\text{d'où}\qquad \operatorname{Ext}^n(k,M)&=\text{annulateur de }\mathfrak m\text{ dans }\operatorname{Ext}^n(A/\mathfrak n,M)
\end{aligned}
\]
donc f) $\Rightarrow$ e), et e) $\Rightarrow H^{n}(M)\subset D(A)$.

% L21–26, lines across earlier argument
\deleted{Donc \unsure{supposons}, si} \inserted{Or l'annulateur} \deleted{$J$ est l'annulateur de $H^n(M)$, on a $H^n(M)\simeq D(A')$ $A'=A/J$. Or $J$ est l'intersection des idéaux premiers de corang $n$ associés à $M$. Or $H^n(M)$ \deleted{\unsure{ne}} change pas si on remplace $M$ par \deleted{\illegible{1}} $M/M_1$, où $M_1$ est l'intersection des noyaux des $M\to M_{\mathfrak p}$ ($\mathfrak p$ de corang $n$).} de $H^n(M)$ est celui de $M$ ($M$ Coh. Mac. donc \uline{équidimensionnel}) donc $M$ étant fidèle, il est nul ; cela signifie (supposons $A$ complet, ce qui est loisible) que $H^n(M)=D(A)$, i.e. b').

% L27–31
\deleted{\unsure{Dans cette question|Donc}} Il reste donc essentiellement les

\uline{Problème} A Si \inserted{dans} $M/x_1M+\ldots+x_nM$, l'annulateur de $\mathfrak m$ est de dim $1$, et \inserted{$(\dim M=n)$} prouver que dans $M$ il est nul

\uline{Problème} B Si \inserted{dans} $\operatorname{Ext}^n(A/\mathfrak n,M)$, \unsure{le annulateur|l'annulateur} de $\mathfrak m$ est de dim $1$, et \inserted{$(\dim M=\dim A=n)$} prouver que dans $M$ \inserted{et $A$} il est nul

% archive sheet number: 36
